% 附录三 颜珰关于中国礼仪布告
% 原书页码：344；标题
\BookUnnumbered{附录三\quad 颜珰关于中国礼仪布告}
\par
% 原书页码：344；段落
节选自《中国礼仪之争：西方文献100篇（1645—1911）》，文献6。%
\par
% 原书页码：344；段落
整篇文献是罗马对颜珰1693年关于中国礼仪的布告的评估。整个%
评估包括四个部分：（1）布告本文；（2）引发的问题；（3）圣职部对这%
些问题的答复；（4）克莱孟十一世批准圣职部的答复的谕令。本附录%
选取的是文献中颜珰的布告部分。%
% 原书页码：344；标题
\BookSubUnnumbered{I. 福建宗座代牧、尊敬的颜珰主教的训令}
\par
% 原书页码：344；段落
{\bookkai 颜珰，司铎，福建的宗座代牧，巴黎神学院博士，索邦神学院院}%
{\bookkai 士，在基督教内向本代牧区全体司铎致意，基督是众人的真正的救主。}%
\par
% 原书页码：344–345；段落（跨页续段）
{\bookkai 自从我们在天主的眷顾之下来到中国之日起，我们的宗座代牧区}%
{\bookkai 所一直关注的问题是在可能的情况下，尽快结束传教士们之间长期以}%
{\bookkai 来争论的种种问题。另外在圣座解决纷争之前，要给予传教士一个准}%
{\bookkai 则，让他们在理论和实践上都可以遵守。我们很遗憾地注意到，在有}%
{\bookkai 关崇敬天主和铲除偶像崇拜的事情上，传播福音的使者们之间意见不}%
{\bookkai 一。这些不同意见经常给教堂带来严重损失。有的人认为一些行为犯}%
{\bookkai 了偶像崇拜罪，而另外一些人则认为不然，允许基督徒去做，或劝他}%
{\bookkai 们去做。尽管有些人认为某些活动是迷信的，但为了求得安宁，还是}%
{\bookkai 暂时允许了这些活动。长期的经验告诉我们，对邪恶的宽容，反而使}%
{\bookkai 得邪恶得以生根并一天天地蔓延。}%
\par
% 原书页码：345；段落
{\bookkai 1. 我们教区的不少传教士在如此重大的问题上，不仅非常希望消}%
{\bookkai 除不同意见，而且还要在行动上保持一致。出于对我们这些管理者的}%
{\bookkai 信任，他们征求我们的意见。他们不愿意始终处于无休止的忐忑不安}%
{\bookkai 中，他们坚持要求我们暂时做出决定。由于这问题比较严重，我们没}%
{\bookkai 有如他们所希望的那样很快给予答复。前任福建代牧区的主教把其任}%
{\bookkai 务交给了我们。我们在很长一段时间内，都在按照他的命令全神贯注}%
{\bookkai 地思考这些问题。在教廷选择了我们来管理这个代牧区后，似乎我们}%
{\bookkai 需要更加勤奋努力地来解决这些问题。我们可以千真万确地说，我们}%
{\bookkai 并没有蓄意忽视过任何能帮助我们解决这些问题的东西，不管是出自}%
{\bookkai 中文的资料，还是某一种欧洲语言写下的评论，或者出自于与学者的}%
{\bookkai 谈话。我们在每天的弥撒及热切的祈祷中祈求天主在黑暗中照亮我们，}%
{\bookkai 教给我们走他的道路，告诉我们在他的眼睛里什么是他喜欢的，是他}%
{\bookkai 欢迎的，从而来解决这些问题。}%
\par
% 原书页码：345–346；段落（跨页续段）
{\bookkai 2. 教廷的敕令和宪章明确指出，主教有责任在自已的代牧区范围}%
{\bookkai 内为崇拜天主及维护伦理道德作好准备。为了履行我们的职责，我们}%
{\bookkai 命令我们教区内的每一个传教士在教廷另做出新决定之前，都要遵守}%
{\bookkai 下列规定：}%
\par
% 原书页码：346；段落
{\bookkai I. 除了特殊情况外，不要用中文去表达无法表达的欧洲名词。我}%
{\bookkai 们宣布应称真神为天主（天上的主人），这已经是用了很久的名称了。}%
{\bookkai 另外两个词——天和上帝（最高的皇帝）——应该完全取消用来翻译}%
{\bookkai 真神。不要让任何一个人知道在汉语中——天和上帝——是我们基督}%
{\bookkai 徒崇拜的真神。}%
\par
% 原书页码：346；段落
{\bookkai Ⅱ.我们严禁在任何教堂里置放刻有“敬天”二字的匾。不管在}%
{\bookkai 什么地方，放上这匾的，我们都下令在两个月之内去掉。凡把真神称为}%
{\bookkai 天或上帝的其他的牌子和类似意义的对联也都要去掉。我们的意见是，}%
{\bookkai 所有这些牌子，尤其那些有“敬天”字样的不能被认为与偶像崇拜无}%
{\bookkai 关。即使事情并没有我们看得那么严重，但也真正暗示着有危险。经验}%
{\bookkai 告诉我们，我们代牧区内那些不让把这些牌子放在自己教堂里的传教}%
{\bookkai 士们，在传播福音方面毫不逊色，他们的收获并不比其他传教士们小。}%
\par
% 原书页码：346；段落
{\bookkai Ⅲ.我们声明：我们之前呈交给教宗亚历山大七世\AuthorNote{〔教皇亚历山大}}%
{\bookkai \AuthorNote{七世之前同意中国礼仪的施行〕}的文件，没有如实反映许多问题，从}%
{\bookkai 而教宗认可了中国人中间已经盛行的祭孔和祭祖。因此，传教士们不}%
{\bookkai 能继续依赖教宗之前做出的答复——教宗的答复当然是正确，明智的，}%
{\bookkai 但却是根据可疑及含糊地描绘的情况而做出的。}%
\par
% 原书页码：346；段落
{\bookkai IV. 传教士在任何情况下都不能允许基督徒主持、参与或者出席}%
{\bookkai 一年数度例行的祭孔、祭祖的隆重仪式，我们宣布这种供祭是带有迷}%
{\bookkai 信色彩。}%
\par
% 原书页码：347；段落
{\bookkai V. 有的传教士在他们传播福音的地方力求取消信徒在家里供先}%
{\bookkai 人的牌位，我们对这些传教士大加赞扬，我们鼓励他们继续做下去。}%
{\bookkai 但是，有的地方很难做到这一点。在这些地方，至少采取这种方法：}%
{\bookkai 把“神主”“神位”“灵位”等字取掉，只能把亡者的名字写在牌位上，}%
{\bookkai 最多只能加一个“位”字。在教廷对这些问题做出裁决以前，我们可}%
{\bookkai 以接受上述形式出现的此种牌位。但是对这种牌位，我们也不能同意}%
{\bookkai 以迷信的眼光来看待。所以，在家里，牌位放置的地方必须用粗体字}%
{\bookkai 写上一个声明，声明基督徒对死者的想法和子孙应该对祖宗行孝道。}%
{\bookkai 我们在这个指令的结尾附了此种声明的一个具体实例。我们并不反对}%
{\bookkai 采用其他类似意义的声明，不过它首先应得到我们的同意。}%
\par
% 原书页码：347；段落
{\bookkai VI. 我们注意到有些在口头上，或在书面上流传着的一些说法，}%
{\bookkai 正把单纯的人引向错误，甚至为迷信打开方便之门。}%
\par
% 原书页码：347；段落
{\bookkai 这些说法包括：}%
\par
% 原书页码：347；段落
{\bookkai a. 中国人教授的哲学，如正确加以理解并没有什么和基督教教}%
\par
% 原书页码：347；段落
{\bookkai 规相违背的；}%
\par
% 原书页码：347；段落
{\bookkai b. 古代的贤人用“太极”这词，把天主解释为世界一切事物的}%
\par
% 原书页码：347；段落
{\bookkai 缘由；}%
\par
% 原书页码：347；段落
{\bookkai c. 孔子向神灵所致的敬意其世俗意义更甚于宗教意义；}%
\par
% 原书页码：347；段落
{\bookkai d. 中国的《礼经》（有关中国所有的习俗和礼仪的书）无论对自}%
\par
% 原书页码：347；段落
{\bookkai 然界还是道德修养来说，都是一部集大成之作。}%
\par
% 原书页码：347；段落
{\bookkai 在我们整个代牧区，我们严格禁止散播大量这类似是而非、错误}%
{\bookkai 百出的言论或文字。}%
\par
% 原书页码：347–348；段落（跨页续段）
{\bookkai VⅦ. 在学校里用中文讲课的基督徒们，不向他们的学生灌输这些}%
{\bookkai 书的原文及注释中层出不穷的无神论及其他迷信的东西。传教士们应}%
{\bookkai 该要求学生驳斥明显的错误，利用批驳错误的机会，趁机孜孜不倦地}%
{\bookkai 教导学生，基督教是如何解释天主及其创造世界、主宰世界的。传教}%
{\bookkai 士们应经常告诫这些基督徒们不要把他们在课外学到的违背基督教义}%
{\bookkai 的东西掺杂进他们的著作中，这种情况是很容易发生的。}%
\par
% 原书页码：348；段落
{\bookkai 其他有些问题可能我们目前还没有考虑到。上述规定是在可靠的}%
{\bookkai 信息来源的基础上制订的。这些信息来源在很大范围内也是其他一些}%
{\bookkai 事情的根据。这些规定指出了实践中应该如何去做。}%
\par
% 原书页码：348；段落
{\bookkai 如果某个传教士——但愿此事不会发生——没有按我们的指令在}%
{\bookkai 两个月内把他管辖的教堂里或屋子里的牌位和我们所指出的条幅拿掉}%
{\bookkai 的话，或者在规定的时间内不努力把我们的准则带进当地的习俗，我}%
{\bookkai 们将撤销我们及任何代牧发放给他的权力，并且我们还要在某个指定}%
{\bookkai 的时间当众宣布。}%
\par
% 原书页码：348；段落
{\bookkai 我们现在重申，并且公开声明，我们并不想给那些以前有另外的}%
{\bookkai 想法，但又不按我们命令去行事的人找岔子。所有传教士们如果意见}%
{\bookkai 不一致，各行其是，每人都认为自己的想法和做法更与真理一致，这}%
{\bookkai 是不足为怪的。但现在让我们大家尝试达到一个共识，大家统一意见，}%
{\bookkai 这样我们将会证明教会并没有被带上污点或缺点，它还是它，神圣而}%
{\bookkai 又纯洁。和平与博爱的天主和我们在一起。}%
\par
% 原书页码：348；段落
{\bookkai 1693年3月26日于福建省长乐}%
\par
% 原书页码：348；段落
{\bookkai 颜珰}%
\par
% 原书页码：348；段落
{\bookkai 福建宗座代牧}%
